\documentclass[11pt]{article}
\usepackage[margin=25mm]{geometry}
\usepackage{fontspec}
\setmainfont{Latin Modern Roman}
\usepackage{amsmath,amssymb,mathtools}
\usepackage{unicode-math}
\setmathfont{Latin Modern Math}
\usepackage{xurl,hyperref}
\hypersetup{colorlinks=true,urlcolor=blue}
\setcounter{secnumdepth}{0}
\setlength{\parindent}{0pt}
\setlength{\parskip}{5pt}
\setlength{\emergencystretch}{3em}
\newcommand{\sref}[2]{\hyperlink{source:#1:#2}{[S#2]}}
\newcommand{\cataloguescope}[1]{\paragraph{Recorded target.}#1}
\begin{document}
\section{Odd-cycle transversal in triangle-free graphs}\label{statement}

\subsection{Definitions and mathematical statement}

\paragraph{Shared notation.}$\mathbb N=\{1,2,\ldots\}$, and $\mathbb Z,\mathbb Q,\mathbb R,\mathbb C$ denote the integers, rationals, reals and complex numbers. Any other conventions are specified locally.


Let $G=(V,E)$ be a finite simple graph, $n=|V|$, and suppose $G$ has no triangle. A set $F\subseteq E$ is an odd-cycle edge transversal if every cycle of odd length contains an edge of $F$; equivalently, $G-F$ is bipartite.
\paragraph{Statement recorded in the supplied source.}
Does every such $G$ have an odd-cycle edge transversal $F$ satisfying $|F|\leq n^2/25$? Equivalently, does $G$ admit a vertex bipartition having at most $n^2/25$ edges within its parts? \sref{statement}{2}
\subsection{Short English statement}
Can every triangle-free graph be made bipartite by deleting at most one twenty-fifth of the square of its number of vertices?
\subsection{Sources}
\begin{description}
\item[\hypertarget{source:statement:1}{S1}] ULAM.AI and the UnsolvedMath Contributors, UnsolvedMath: A Curated Collection of Open Mathematics Problems, catalogue entry. CC BY 4.0. \url{https://huggingface.co/datasets/ulamai/UnsolvedMath}.
\item[\hypertarget{source:statement:2}{S2}] Open Problem Garden, Odd-cycle transversal in triangle-free graphs, node 46824, original problem page. Original question and full discussion. \url{https://www.openproblemgarden.org/op/odd_cycle_transversal_in_triangle_free_graphs}.

\end{description}
\label{end:statement}
\end{document}
